% LuaLaTeX document — compile with: lualatex mie_multipole_derivation.tex
\documentclass[11pt,a4paper]{article}

\usepackage{fontspec}
\setmainfont{Latin Modern Roman}
\setmonofont{Latin Modern Mono}

\usepackage{amsmath,amssymb,amsthm,bm}
\usepackage{booktabs}
\usepackage{geometry}
\geometry{margin=2.5cm}
\usepackage{hyperref}
\hypersetup{colorlinks=true,linkcolor=blue!60!black,citecolor=green!50!black,urlcolor=blue!70!black}
\usepackage{xcolor}
\usepackage{mathtools}
\usepackage{enumitem}

\newcommand{\D}{\Delta}
\newcommand{\order}{\mathcal{O}}
\newcommand{\code}[1]{\texttt{#1}}
\newcommand{\bu}{\mathbf{u}}
\newcommand{\br}{\mathbf{r}}
\newcommand{\bM}{\mathbf{M}}
\newcommand{\bP}{\mathbf{P}}
\newcommand{\bS}{\mathbf{S}}
\newcommand{\bnabla}{\boldsymbol{\nabla}}
\newcommand{\bpsi}{\boldsymbol{\psi}}
\newcommand{\bphi}{\boldsymbol{\varphi}}
\newcommand{\hatphi}{\hat{\mathbf{e}}_\varphi}

\title{Direct Derivation of the Elastic Mie Multipole Series\\
\large From Helmholtz potentials to closed-form scattering amplitudes}
\author{T.~M.~Nestor}
\date{April 2026}

\begin{document}
\maketitle

\begin{abstract}
We give a self-contained derivation of the elastic Mie partial-wave
series for a spherical inclusion in an isotropic background, ending at
the closed-form scattering amplitudes $a_n,b_n,c_n,d_n$ used by
\code{Mathematica/MieAsymptotic.wl}.  The derivation is organised in
eight steps: Helmholtz decomposition, Rayleigh expansion of the incident
plane wave, three-region partial-wave ansatz, field-from-potential
operators, boundary conditions at $r=a$, Cramer's rule, series expansion
in $ka$, and the appearance of the static Eshelby concentration as a
denominator factor.  The Eshelby concentration falls out of the
determinant arithmetic --- it is never inserted by hand.
\end{abstract}

\tableofcontents

% =====================================================================
\section{Setting and Helmholtz decomposition}
% =====================================================================

In an isotropic linearly elastic medium with Lam\'e parameters
$(\lambda,\mu)$ and density $\rho$, the displacement $\bu(\br)$ at
angular frequency $\omega$ obeys the Navier equation
\begin{equation}
(\lambda+\mu)\,\bnabla(\bnabla\!\cdot\!\bu) + \mu\,\nabla^2\bu + \rho\omega^2\bu \;=\; 0.
\label{eq:navier}
\end{equation}
The two characteristic wavespeeds are
$\alpha=\sqrt{(\lambda+2\mu)/\rho}$ and $\beta=\sqrt{\mu/\rho}$, with
wavenumbers $k_P=\omega/\alpha$ and $k_S=\omega/\beta$.

Any solution of \eqref{eq:navier} can be written as the sum of a
curl-free (compressional) and a divergence-free (shear) part:
\begin{equation}
\bu \;=\; \bnabla\phi \;+\; \bnabla\!\times\!\bpsi.
\label{eq:helmholtz-decomp}
\end{equation}
For a $P$-wave incident along $\hat{\mathbf z}$, the problem is
axisymmetric ($\partial_\varphi\equiv 0$) and the vector potential
reduces to a single component, $\bpsi=\psi(r,\theta)\,\hatphi$.  The
scalar $\phi$ and the scalar $\psi$ each satisfy a scalar Helmholtz
equation:
\begin{equation}
(\nabla^2 + k_P^2)\,\phi = 0, \qquad
(\nabla^2 + k_S^2 - r^{-2}\sin^{-2}\theta)\,\psi = 0,
\label{eq:helmholtz-scalar}
\end{equation}
the second arising from the vector Laplacian acting on
$\psi\,\hatphi$.  Substituting back into \eqref{eq:helmholtz-decomp}
gives the explicit components
\begin{equation}
u_r = \partial_r\phi + \frac{1}{r\sin\theta}\,\partial_\theta(\sin\theta\,\psi),
\qquad
u_\theta = \frac{1}{r}\partial_\theta\phi - \frac{1}{r}\partial_r(r\,\psi).
\label{eq:u-from-pot}
\end{equation}

% =====================================================================
\section{Rayleigh expansion of the incident plane wave}
\label{sec:rayleigh}
% =====================================================================

The incident $P$-wave $\bu^{\rm inc}=\bnabla\phi^{\rm inc}$ with unit
amplitude scalar potential
$\phi^{\rm inc}(\br)=e^{ik_P z}/(ik_P)$ admits the
Gegenbauer/Rayleigh expansion in spherical harmonics:
\begin{equation}
\phi^{\rm inc}(r,\theta) \;=\; \frac{1}{ik_P}
\sum_{n=0}^{\infty}(2n+1)\,i^n\,j_n(k_P r)\,P_n(\cos\theta),
\label{eq:rayleigh-exp}
\end{equation}
where $j_n$ is the spherical Bessel function of the first kind.  The
factor $1/(ik_P)$ in front is conventional: it makes
$\partial_z e^{ik_P z}/(ik_P) = e^{ik_P z}$, so the displacement
amplitude has unit normalisation along the axis.  Equation
\eqref{eq:rayleigh-exp} decouples the problem into a direct sum over
Legendre orders $n$: each $n$ propagates independently through the
spherically symmetric inclusion.

The companion expansion for an incident $SV$-wave uses the associated
Legendre polynomials $P_n^1(\cos\theta)$ with the vector-potential
form; we restrict here to $P$-incidence.

% =====================================================================
\section{Three-region partial-wave ansatz}
\label{sec:ansatz}
% =====================================================================

Inside the inclusion ($r<a$) the wavenumbers are primed:
$k_P'=\omega/\alpha'$, $k_S'=\omega/\beta'$ with
$\alpha'=\sqrt{(\lambda+\D\lambda+2(\mu+\D\mu))/(\rho+\D\rho)}$ etc.
Outside ($r>a$) the wavenumbers are unprimed.  At each Legendre order
$n$ the potentials take the form
\begin{equation}
\begin{aligned}
\phi^{\rm inc}_n &= \frac{(2n+1)i^n}{ik_P}\,j_n(k_P r)\,P_n(\cos\theta), \\[4pt]
\phi^{\rm sc}_n &= a_n\,h_n^{(1)}(k_P r)\,P_n, &
\psi^{\rm sc}_n &= b_n\,h_n^{(1)}(k_S r)\,P_n^1, \\[4pt]
\phi^{\rm int}_n &= c_n\,j_n(k_P' r)\,P_n, &
\psi^{\rm int}_n &= d_n\,j_n(k_S' r)\,P_n^1,
\end{aligned}
\label{eq:ansatz}
\end{equation}
where:
\begin{itemize}[leftmargin=*,nosep]
\item $j_n$ (regular at the origin) is the only physically admissible
radial function inside the inclusion;
\item $h_n^{(1)}=j_n+iy_n$ (Hankel, outgoing at infinity) is the only
physically admissible scattered radial function outside.
\end{itemize}
For $n=0$ the associated Legendre $P_0^1\equiv 0$, so the
$\psi$-channel is empty: only $a_0$ and $c_0$ remain, and the
boundary system is $2\times 2$.  For $n\ge 1$ all four amplitudes
$(a_n,b_n,c_n,d_n)$ are non-trivial.

% =====================================================================
\section{Displacement and stress from a partial-wave potential}
\label{sec:fields}
% =====================================================================

Write $z_n(\xi)$ for a generic spherical Bessel/Hankel of order $n$
($z_n\in\{j_n,h_n^{(1)}\}$).  Substituting a single partial-wave
potential
$\phi=z_n(kr)\,P_n(\cos\theta)$ or $\psi=z_n(kr)\,P_n^1(\cos\theta)$
into Eq.~\eqref{eq:u-from-pot} and applying Hooke's law
\begin{equation}
\sigma_{ij} = \lambda\,\delta_{ij}\,\varepsilon_{kk} + 2\mu\,\varepsilon_{ij},
\qquad
\varepsilon_{ij} = \tfrac{1}{2}(\partial_i u_j + \partial_j u_i),
\label{eq:hooke}
\end{equation}
gives the four boundary quantities $u_r$, $u_\theta$, $\sigma_{rr}$,
$\sigma_{r\theta}$ in compact form.  With shorthand $\xi\equiv kr$:

\paragraph{P-channel ($\phi$-potential).}
\begin{equation}
\begin{aligned}
u_r &= k\,z_n'(\xi)\,P_n, \\
u_\theta &= \frac{z_n(\xi)}{r}\,\partial_\theta P_n, \\
\sigma_{rr} &= -(\lambda+2\mu)\,k^2\,z_n(\xi)\,P_n
              \;-\;\tfrac{4\mu k}{r}z_n'(\xi)\,P_n
              \;+\;\tfrac{2\mu n(n{+}1)}{r^2}\,z_n(\xi)\,P_n, \\
\sigma_{r\theta} &= \tfrac{2\mu}{r}\Bigl[k\,z_n'(\xi) - \tfrac{z_n(\xi)}{r}\Bigr]\,\partial_\theta P_n.
\end{aligned}
\label{eq:p-fields}
\end{equation}

\paragraph{S-channel ($\psi$-potential).}
\begin{equation}
\begin{aligned}
u_r &= \frac{n(n{+}1)\,z_n(\xi)}{r}\,P_n, \\
u_\theta &= \frac{z_n(\xi) + \xi z_n'(\xi)}{r}\,\partial_\theta P_n, \\
\sigma_{rr} &= \tfrac{2\mu n(n{+}1)}{r}\Bigl[k\,z_n'(\xi) - \tfrac{z_n(\xi)}{r}\Bigr]\,P_n, \\
\sigma_{r\theta} &= \tfrac{\mu}{r^2}\bigl[(2n(n{+}1)-2-\xi^2)\,z_n(\xi) - 2\xi\,z_n'(\xi)\bigr]\,\partial_\theta P_n.
\end{aligned}
\label{eq:s-fields}
\end{equation}
These four-vectors are exactly the columns assembled by the
\code{pField} and \code{sField} functions in
\code{Mathematica/MieAsymptotic.wl}.

% =====================================================================
\section{Boundary conditions: the $4\times 4$ system}
\label{sec:bc}
% =====================================================================

At the spherical interface $r=a$ the elastic boundary conditions are
continuity of (i) the displacement vector and (ii) the traction vector
$t_i=\sigma_{ij}\hat n_j$.  Since $\hat n=\hat r$, this means
continuity of the four scalar quantities
$\{u_r,u_\theta,\sigma_{rr},\sigma_{r\theta}\}$.  Imposing these on
the sum (incident + scattered) outside and the interior expansion
inside gives, per partial wave $n\ge 1$,
\begin{equation}
\bM_n(w;\,\D\lambda,\D\mu,\D\rho)\,
\begin{pmatrix} a_n \\ b_n \\ c_n \\ d_n \end{pmatrix}
\;=\; -\,\frac{(2n+1)\,i^n}{ik_P}\,\mathbf{f}^{(P,j)}_n(a),
\label{eq:bc-system}
\end{equation}
where $w\equiv k_S a$ is the dimensionless small parameter and the
$4\times 4$ matrix has columns
\begin{equation}
\bM_n \;=\; \bigl[\;\mathbf{f}^{(P,h)}_n(a)\;\;\;\mathbf{f}^{(S,h)}_n(a)\;\;\;
                 -\mathbf{f}^{(P,j')}_n(a)\;\;\;-\mathbf{f}^{(S,j')}_n(a)\;\bigr].
\label{eq:bc-matrix}
\end{equation}
Here $\mathbf{f}^{(\bullet)}_n(a)\in\mathbb{C}^4$ collects the four
boundary quantities of \eqref{eq:p-fields}--\eqref{eq:s-fields} at
$r=a$, evaluated with the appropriate radial function ($j_n$ or
$h_n^{(1)}$) and the appropriate moduli (background or interior).
The right-hand side is the contribution of the incident wave evaluated
at $r=a$ with $z_n=j_n$ and the background moduli.  For $n=0$ the
matrix is $2\times 2$ (only $a_0,c_0$ unknown, only $u_r,\sigma_{rr}$
constraints).

\subsection*{Why this is a closed problem}
We have four unknowns $(a_n,b_n,c_n,d_n)$ and four equations from
\eqref{eq:bc-system}.  All matrix entries are exact functions of the
contrasts $(\D\lambda,\D\mu,\D\rho)$ and of $w$; no approximation has
been made.

% =====================================================================
\section{Cramer's rule}
\label{sec:cramer}
% =====================================================================

Each amplitude is a ratio of $4\times 4$ determinants:
\begin{equation}
a_n \;=\; \frac{\det \bM_n^{[1\to\mathbf{r}]}}{\det \bM_n},
\qquad
b_n \;=\; \frac{\det \bM_n^{[2\to\mathbf{r}]}}{\det \bM_n},
\qquad \dots
\label{eq:cramer}
\end{equation}
where $\bM_n^{[i\to\mathbf{r}]}$ denotes $\bM_n$ with column $i$
replaced by the right-hand side of \eqref{eq:bc-system}.  This is the
\emph{exact} closed form: at this stage every amplitude is a rational
function of $w$ and the contrasts.  No series expansion has been taken.

A practical complication is that $\det\bM_n$ has a non-trivial leading
order in $w$ (it scales as $w^{4n+2}$ for the $4\times 4$ block, modulo
the convention factor in the RHS), and $\det\bM_n^{[1\to\mathbf{r}]}$
scales differently --- so the ratio is a genuine \emph{Laurent}
expansion in $w$, not a pure Taylor series.  Both numerator and
denominator must therefore be expanded together and the leading
\emph{ratio} extracted.

% =====================================================================
\section{Series expansion in $ka$}
\label{sec:series}
% =====================================================================

The Bessel and Hankel functions admit known power-series expansions
near $\xi=0$:
\begin{equation}
j_n(\xi) \;=\; \frac{\xi^n}{(2n+1)!!}\biggl[1-\frac{\xi^2}{2(2n+3)}+\order(\xi^4)\biggr],
\qquad
y_n(\xi) \;=\; -\frac{(2n-1)!!}{\xi^{n+1}}\bigl[1+\order(\xi^2)\bigr].
\label{eq:bessel-series}
\end{equation}
Substituting into \eqref{eq:cramer} via the Sin/Cos recurrence
implementation of \code{MieAsymptotic.wl} (lines defining
\code{jn}, \code{yn}, \code{hn}) makes both numerator and denominator
of $a_n$ Laurent polynomials in $w$ with coefficients that are
\emph{exact rational functions in $(\D\lambda,\D\mu,\D\rho,\lambda_0)$}.
Series-dividing the ratio (Mathematica's \code{Series[\dots,\{w,0,\text{ord}\}]}
applied to \code{Det[col]/dM}) collapses the Laurent expansion and
extracts the leading coefficient.

Carrying out this procedure analytically for $n=0,1,2$ produces
\begin{align}
a_0 &= \frac{w^2(-3\D\lambda-2\D\mu)}{3(\lambda_0+2)(6+3\D\lambda+2\D\mu+3\lambda_0)}, \\[4pt]
a_1 &= \frac{i\,w^2\,\D\rho}{3(\lambda_0+2)}, \qquad
b_1 \;=\; \frac{i\,w^2\,\D\rho}{3}, \\[4pt]
a_2 &= \frac{20\,w^2\,\D\mu}{3(\lambda_0+2)\bigl[15(\lambda_0+2)+2\D\mu(8+3\lambda_0)\bigr]},
\label{eq:a2-leading}\\[4pt]
b_2 &= \frac{10\,w^2\,\D\mu\,\sqrt{\lambda_0+2}}{3\bigl[15(\lambda_0+2)+2\D\mu(8+3\lambda_0)\bigr]}.
\end{align}
These are the closed forms tabulated in
\code{Mathematica/MieAsymptoticLeading.wl} and exposed to Python via
\code{cubic\_scattering/mie\_asymptotic\_analytic.py}.

\subsection*{Higher-order corrections}
Continuing the series expansion to order $w^{2n+3}$ produces
$\order((ka)^2)$ corrections to each leading amplitude; these encode
the dynamic (radiative) corrections to the static-Eshelby leading
behaviour.  At $ka_S=0.05$ the relative size of the first correction
is $\sim 2.5\times 10^{-3}$, consistent with the Eshelby-vs-Mie
agreement in the validation document.

% =====================================================================
\section{Where the Eshelby concentration comes from}
\label{sec:eshelby-emerges}
% =====================================================================

The denominator factor in $a_2$ of Eq.~\eqref{eq:a2-leading} can be
rewritten as
\begin{equation}
15(\lambda_0+2)+2\D\mu(8+3\lambda_0) \;=\; 15(\lambda_0+2)\,(1+\beta_E\D\mu),
\qquad
\beta_E \;=\; \frac{2(8+3\lambda_0)}{15(\lambda_0+2)}.
\label{eq:eshelby-emerges}
\end{equation}
Comparing with $\beta_E = 2(4-5\nu)/(15(1-\nu))$ via $\nu=\lambda_0/(2(\lambda_0+1))$
shows this is exactly the static Eshelby shear concentration
coefficient for a sphere.  In other words,
\begin{equation}
a_2 \;=\; \frac{4\,w^2\,\D\mu}{9(\lambda_0+2)^2}\;\cdot\;\underbrace{\frac{1}{1+\beta_E\D\mu}}_{E_2,\;\text{Eshelby}}.
\label{eq:a2-factored}
\end{equation}
The Eshelby concentration emerges \emph{automatically} from the
determinant arithmetic of \eqref{eq:cramer}.  Two observations
illuminate why:

\begin{enumerate}[leftmargin=*]
\item\emph{The Cramer ratio is the inverse of an operator equation.}
The Lippmann--Schwinger volume integral equation reads
$\bu^{\rm tot} = \bu^{\rm inc} + (\hat G_0 \hat T_0)\,\bu^{\rm tot}$,
so the polarisation satisfies
$\bP = (I-\hat G_0\hat T_0)^{-1}\,\hat T_0\,\bu^{\rm inc}$.  Inverting
$(I-\hat G_0\hat T_0)$ generates a geometric series in
$\hat G_0\hat T_0$ whose projection onto each angular sector is
exactly the geometric series $1 + \beta_E\D\mu + (\beta_E\D\mu)^2+\dots$
for the shear sector.  The $4\times 4$ Mie BC matrix $\bM_n$ is the
spherical-harmonic projection of this same operator $(I-\hat G_0\hat T_0)$;
its determinant ratio reproduces the geometric series sum
$1/(1+\beta_E\D\mu)$.

\item\emph{The spherical harmonic basis is the diagonal basis of the
Eshelby operator for a sphere.}  Each $n$ is a separate eigenvector of
the static Eshelby tensor.  The Mie BC system thus diagonalises the
Lippmann--Schwinger inversion exactly --- one $4\times 4$ block per
$n$, no coupling between Legendre orders.  This is why the closed-form
$a_n,b_n$ encode finite-contrast Eshelby \emph{without any
approximation}: the polynomial-Galerkin closure on the cube
(\code{cube\_galerkin27.tex}) does the analogous inversion on a
non-diagonal but finite-dimensional basis ($\mathcal{T}_{27}$).
\end{enumerate}

% =====================================================================
\section*{Summary}
\addcontentsline{toc}{section}{Summary}
% =====================================================================

The eight-step recipe --- Helmholtz decomposition, Rayleigh expansion,
three-region ansatz, field operators, BC system, Cramer's rule, series
expansion, Eshelby identification --- is mechanised verbatim in
\code{Mathematica/MieAsymptotic.wl}.  The pieces map as follows:

\begin{center}
\begin{tabular}{ll}
\toprule
Mathematical step & Implementation \\
\midrule
Helmholtz potentials & implicit (the $\phi/\psi$ ansatz of \S\ref{sec:ansatz}) \\
Rayleigh expansion of incident & \code{coeff = (2n+1) I\^{}n/(I kPo)} \\
P-channel field operators
   \eqref{eq:p-fields} & \code{pField[n,k,r,lam,mu,zfn]} \\
S-channel field operators
   \eqref{eq:s-fields} & \code{sField[n,k,r,mu,zfn]} \\
BC matrix \eqref{eq:bc-matrix} & \code{mat = Transpose[\{scP,scS,-inP,-inS\}]} \\
Cramer ratio \eqref{eq:cramer} & \code{Det[col]/dM} per column \\
Series in $w$ \eqref{eq:bessel-series} & \code{Normal[Series[\dots,\{w,0,ord\}]]} \\
Eshelby identification \eqref{eq:eshelby-emerges} &
   automatic --- see \code{MieAsymptoticVerify.wl} \\
\bottomrule
\end{tabular}
\end{center}

\noindent For the validation that the closed forms above reproduce the
spherical Eshelby concentration factors and match the project's
$27\times 27$ cubic T-matrix at finite contrast, see the companion
document \code{mie\_finite\_contrast\_validation.tex}.

\end{document}
